% =====================================================================
\lezione{20/10/2026}{Imparare gli embedding; dropout e batch norm nel dettaglio}
  {Prof.~Danilo Giordano}
  {embedding supervisionati con \code{nn.Embedding}: dati tabellari, sequenze, architetture composte;
   shortcut learning; embedding self-supervised (autoencoder, contrastive learning); dropout e batch
   normalization in training e in valutazione; linee guida per i report dei lab}
\label{lez:2026-10-20}
% =====================================================================

\begin{nota}[Fonte]
Trascrizione automatica della registrazione, revisionata, confrontata con le slide
\emph{4-embedding-introduction} (pp.~23--51), \emph{2a-dropout\_batchnorm} e con il documento di linee guida
per i report presentato a fine lezione.
\end{nota}

\section{Sintesi}
La lezione completa il discorso sulle feature categoriche iniziato il giorno prima (\rif{2026-10-19}{19/10}):
come si \textbf{impara} un embedding. Prima nel caso supervisionato, con tre esempi di complessità
crescente (tabella, sequenza, grafi nel tempo), poi nel caso \textbf{self-supervised}, dove il segnale non
viene dalle label ma dalle proprietà interne dei dati. Seguono un breve deck che chiarisce come
\textbf{dropout} e \textbf{batch normalization} si comportano in training e in valutazione, e la
presentazione delle \textbf{linee guida per i report}. Approfondimenti nei capitoli~\ref{cap:embedding},
\ref{cap:regolarizzazione} e~\ref{cap:lab-report}.

\section{Avvisi pratici}
\begin{itemize}
  \item Il Lab~2 inizia \textbf{il giorno dopo} (mercoledì 21/10): gli embedding visti oggi servono lì.
  \item Sul portale è caricato il documento con le \textbf{raccomandazioni per i report} (dettagli nel
    capitolo~\ref{cap:lab-report}). Punti richiamati in aula:
    \begin{itemize}
      \item chi non è ancora in un gruppo deve unirsi a uno \textbf{entro fine mese}; ci sono gruppi da due o
        tre persone da contattare;
      \item consegna \textbf{10 giorni prima} dello scritto; PDF di \textbf{massimo 10 pagine} con lab ID,
        group ID, nomi ed email nell'intestazione; LaTeX o Word indifferentemente, con un font leggibile;
      \item il notebook deve \textbf{girare ed essere già eseguito} (con gli output visibili), ma il report
        deve essere autosufficiente: non si va a cercare i risultati nel notebook;
      \item novità di quest'anno: un \textbf{logbook degli esperimenti} sintetico, che \textbf{non conta}
        nelle 10 pagine;
      \item ogni report vale da 0 a 32 e serve almeno 18; si valuta il ragionamento, non lo 0,81 contro lo
        0,82; figure leggibili e tabelle al posto di elenchi di numeri nel testo.
    \end{itemize}
\end{itemize}

\section{Contenuti}

\subsection{Embedding supervisionati}
Ipotesi: label disponibili, feature numeriche e categoriche, e una \textbf{rete neurale} (un random forest
classifica ma non fa feature extraction). \code{nn.Embedding(V, H)} è una lookup table addestrabile con $V$
righe (le categorie) e $H$ colonne (dimensione dell'embedding, un iperparametro); la backpropagation ne
aggiorna i valori come per gli altri pesi.
\begin{itemize}
  \item \textbf{Tabella} (metadati di email, task phishing/normale): il dominio del mittente ha cardinalità
    enorme e nessun ordine $\Rightarrow$ embedding. Si fa lo split, si costruisce il \textbf{vocabolario
    solo sul training} (validation e test simulano il futuro) con \code{<UNK>} per i valori mai visti, si
    mappano gli ID, si separano categoriche e numeriche, e nel modello l'embedding si \textbf{concatena}
    alle numeriche prima della FFNN (input del primo layer: \code{embed\_dim + num\_numeric}).
  \item \textbf{Sequenza} (IP che contattano una darknet, task: attacco di botnet e quale): embedding su
    ogni elemento, poi LSTM con input di dimensione \code{embed\_dim}; con \code{padding\_idx} l'embedding
    del padding non viene aggiornato. È lo schema del lab.
  \item \textbf{Architettura composta} (movimento laterale): embedding dei servizi $\to$ GNN per ogni
    snapshot $\to$ RNN sulla serie di snapshot $\to$ predizione, tutto addestrato con lo stesso segnale
    supervisionato.
\end{itemize}

\begin{esame}[un embedding supervisionato vale per il suo task]
Il docente insiste: l'embedding appreso con le label è \textbf{ottimo per quel task e sub-ottimo per
altri}. Numericamente lo si può riusare, ma è stato costruito per un altro obiettivo. E può imparare
\textbf{scorciatoie}: nell'esempio husky contro lupo il classificatore impara lo sfondo (erba o neve),
non l'animale.
\end{esame}

\subsection{Embedding self-supervised}
Senza label il segnale viene dalle \textbf{proprietà interne} dei dati. Due strade: un
\textbf{autoencoder}, in cui l'encoder comprime e il decoder ricostruisce (un problema di regressione senza
label; al test si congela l'encoder e si addestra una piccola head con poche label); il
\textbf{contrastive learning}, in cui le versioni aumentate di un'immagine (positivi) devono avere embedding
vicini all'originale (anchor) e un'immagine diversa (negativo) lontano, con la \textbf{contrastive loss}.
Pre-addestramento su grandi corpora (\en{pretext}) e fine-tuning con poche label (\en{downstream}) sono la
base dei successi recenti, e sono preziosi in cybersecurity dove le label scarseggiano. Le augmentation
però iniettano conoscenza, quindi bias (la domanda della slide sullo scoiattolo verde).

\subsection{Dropout e batch normalization nel dettaglio}
\begin{itemize}
  \item \textbf{Dropout}: in training ogni neurone è spento con probabilità $p$, con una maschera nuova per
    ogni campione e ogni mini-batch, e i sopravvissuti sono scalati di $1/(1-p)$; in valutazione è
    l'identità. Ogni neurone impara solo dai campioni in cui era acceso.
  \item \textbf{Batch norm}: in training standardizza ogni feature con media e varianza \textbf{del
    mini-batch}, poi applica $\gamma$ e $\beta$ \textbf{appresi}, e aggiorna media e varianza
    \textbf{mobili} (calcolate, non apprese); in valutazione usa le statistiche mobili congelate, così il
    layer successivo riceve dati con la stessa distribuzione vista in training.
\end{itemize}
Tabelle ed esempio numerico nel capitolo~\ref{cap:regolarizzazione}.

\begin{integrazione}[\code{model.eval()} e \code{torch.no\_grad()}]
È \code{model.eval()} che cambia il comportamento di dropout e batch norm; \code{torch.no\_grad()}
disattiva soltanto il calcolo dei gradienti. In valutazione servono entrambi.
\end{integrazione}

\section{Domande aperte e punti da verificare}
\begin{verifica}
\begin{itemize}
  \item Domanda aperta della slide: un modello addestrato con augmentation che tolgono il colore
    riconoscerebbe come anomala la foto di uno scoiattolo verde? Le slide la lasciano come spunto di
    ragionamento, senza risposta ufficiale.
  \item Formato e lunghezza del logbook degli esperimenti: da controllare nel documento sul portale.
\end{itemize}
\end{verifica}
